% Part:sets-functions-relations
% Chapter: size-of-sets
% Section: non-enumerability-alt

\documentclass[../../../include/open-logic-section]{subfiles}

\begin{document}

\olfileid[fa-IR]{sfr}{siz}{nen-alt}

\olsection{مجموعه‌های \printtoken{S}{nonenumerable}}

\begin{editorial}
  این بخش ناشمارابودنِ $\Bin^\omega$ و $\Pow{\Nat}$ را با استفاده از
  تعریف‌های \olref[enm-alt]{sec} ثابت می‌کند؛ یعنی وجود تناظری دوسویی
  با~$\Nat$ را به‌جای یک تابع پوشا از $\PosInt$ لازم می‌گیرد.
\end{editorial}

\begin{explain}
مجموعهٔ $\Nat$ از اعداد طبیعی نامتناهی است. همچنین به‌وضوح
!!{enumerable} است. اما واقعیت شگفت‌انگیز این است که مجموعه‌های
\emph{!!{nonenumerable}} وجود دارند؛ یعنی مجموعه‌هایی که !!{enumerable}
نیستند (نگاه کنید به \olref[sfr][siz][enm-alt]{defn:enumerable}).

این مطلب شاید شگفت‌آور باشد. زیرا گفتن اینکه $A$
!!{nonenumerable} است یعنی گفتن اینکه \emph{هیچ} !!{bijection}‌ای مانند $f
\colon \Nat \to A$ وجود ندارد؛ یعنی هیچ تابعی که !!{element}sی~$\Nat$
را، که شمارشان نامتناهی است، به~$A$ نگاشت کند، همهٔ~$A$ را فرا نمی‌گیرد. پس اگر
$A$ !!{nonenumerable} باشد، شمار !!{element}sی~$A$ از شمار اعداد طبیعی
«بیشتر» است.

برای اثبات اینکه مجموعه‌ای !!{nonenumerable} است، باید نشان دهید که هیچ
!!{bijection} مناسبی نمی‌تواند وجود داشته باشد. بهترین راه برای این کار
آن است که نشان دهید هر کوششی برای برشمردن !!{element}sی~$A$ ناگزیر
دست‌کم یک !!{element} را از قلم می‌اندازد؛ این نشان می‌دهد که هیچ تابع
$f\colon \Nat \to A$ !!{surjective} نیست. راهبردی کلی برای اثبات این
مطلب، استفاده از \emph{روش قطری} کانتور است. با داشتن فهرستی از
!!{element}sی $A$، مثلاً $x_1$، $x_2$، \dots، !!{element} دیگری از~$A$
می‌سازیم که بنا بر نحوهٔ ساختش، به‌هیچ‌وجه نمی‌تواند در آن فهرست باشد.

اما همهٔ این مطالب با یک مثال بهتر فهمیده می‌شوند. بنابراین نخستین مثال
ما مجموعهٔ~$\Bin^\omega$ از همهٔ رشته‌های نامتناهیِ $0$ها و $1$هاست.
(نماد $\Bin$ نمایندهٔ «دودویی» است و می‌توانیم آن را صرفاً مجموعهٔ دوعضویِ
$\{0,1\}$ بدانیم.)\oliflabeldef{sfr:card-arithmetic:card-opps:sec}{\footnote{دقیق‌تر
آن است که مقرر کنیم $\Bin^\omega$ مجموعهٔ همهٔ دنباله‌های $\omega$تایی
از $0$ها و $1$ها، یعنی مجموعهٔ
$\funfromto{\omega}{\{0,1\}}$، است. اما معنای این مطلب تنها در
\olref[sfr][card-arithmetic][card-opps]{sec} روشن خواهد شد.}{} بااین‌حال،
این صورت‌بندیِ اندکی مسامحه‌آمیز فعلاً نباید موجب هیچ ابهامی شود.}
\end{explain}

\begin{thm}
\ollabel{thm:nonenum-bin-omega}
$\Bin^\omega$~!!{nonenumerable} است.
\end{thm}

\begin{proof}
هر برشماری از یک زیرمجموعهٔ $\Bin^\omega$ را در نظر بگیرید. پس فهرستی
مانند $s_{0}$، $s_{1}$، $s_{2}$، \dots{} داریم که در آن هر $s_n$ رشته‌ای
نامتناهی از $0$ها و~$1$هاست. بگذارید $s_n(m)$ رقم $n$اُمِ رشتهٔ $m$اُمِ
این فهرست باشد. اکنون می‌توانیم فهرست خود را آرایه‌ای در نظر بگیریم که
در آن $s_n(m)$ در سطر $n$اُم و ستون $m$اُم قرار گرفته است:
\[
\begin{array}{c|c|c|c|c|c}
& 0 & 1 & 2 & 3 & \dots \\\hline
0 & \mathbf{s_{0}(0)} & s_{0}(1) & s_{0}(2) & s_0(3) & \dots \\\hline
1 & s_{1}(0)& \mathbf{s_{1}(1)} & s_1(2) & s_1(3) & \dots \\\hline
2 & s_{2}(0)& s_{2}(1) & \mathbf{s_2(2)} & s_2(3) & \dots \\\hline
3 & s_{3}(0)& s_{3}(1) & s_3(2) & \mathbf{s_3(3)} & \dots \\\hline
\vdots & \vdots & \vdots & \vdots & \vdots & \mathbf{\ddots}
\end{array}
\]
اکنون رشته‌ای نامتناهی، یعنی $d$، از $0$ها و $1$ها می‌سازیم که در این
فهرست نیست. این کار را با مشخص‌کردن هریک از درایه‌های آن انجام می‌دهیم؛
یعنی $d(n)$ را برای همهٔ $n \in \Nat$ مشخص می‌کنیم. به‌طور شهودی، قطر
آرایهٔ بالا را رو به پایین می‌خوانیم (ازاین‌رو نام «روش قطری») و سپس هر
$1$ را به $0$ و هر $1$ را به~$0$ تغییر می‌دهیم. به‌بیان انتزاعی‌تر،
$d(n)$ را برابر $0$ یا $1$ تعریف می‌کنیم، بسته به اینکه !!{element}
$n$اُمِ قطر، یعنی $s_n(n)$، برابر $1$ یا $0$ باشد؛ یعنی:
\[
d(n) =
\begin{cases}
1 & \text{اگر $s_{n}(n) = 0$}\\
0 & \text{اگر $s_{n}(n) = 1$}
\end{cases}
\]
آشکارا $d \in \Bin^\omega$ است، زیرا رشته‌ای نامتناهی از $0$ها و $1$هاست.
اما $d$ را چنان ساخته‌ایم که $d(n) \neq s_n(n)$ برای هر $n \in \Nat$.
یعنی $d$ با $s_n$ در درایهٔ $n$اُم تفاوت دارد. پس $d \neq s_n$
برای هر $n\in \Nat$. بنابراین $d$ نمی‌تواند در فهرست
$s_0$، $s_1$، $s_2$،
\dots باشد.

نشان داده‌ایم که هر برشماریِ دلخواه از زیرمجموعه‌ای از $\Bin^\omega$،
یک !!{element} از $\Bin^\omega$ را از قلم می‌اندازد. پس هیچ برشماری‌ای
از مجموعهٔ $\Bin^\omega$ وجود ندارد؛ یعنی $\Bin^\omega$
!!{nonenumerable} است.
\end{proof}

\begin{explain}
این روش اثبات را «قطری‌سازی» می‌نامند، زیرا برای تعریف~$d$ از قطر آرایه
استفاده می‌کند. بااین‌حال، قطری‌سازی لزوماً مستلزم وجود آرایه نیست. در
واقع، می‌توانیم با اندیشه‌ای مشابه نشان دهیم که مجموعه‌ای
!!{nonenumerable} است، حتی هنگامی که هیچ آرایه و هیچ قطر واقعی‌ای در
میان نیست. نتیجهٔ زیر چگونگی این کار را نشان می‌دهد.
\end{explain}

\begin{thm}
\ollabel{thm:nonenum-pownat}
$\Pow{\Nat}$ !!{enumerable} نیست.
\end{thm}

\begin{proof}
به همان شیوه پیش می‌رویم و نشان می‌دهیم که هر فهرستی از زیرمجموعه‌های
~$\Nat$، زیرمجموعه‌ای از $\Nat$ را از قلم می‌اندازد. پس فرض کنید فهرستی
مانند $N_0, N_1, N_2, \ldots$ از زیرمجموعه‌های $\Nat$ داریم. مجموعهٔ
$D$ را چنین تعریف می‌کنیم: $n \in D$ اگر و تنها اگر $n \notin N_{n}$:
\[
D = \Setabs{n \in \Nat}{n \notin N_n}
\]
آشکارا $D\subseteq \Nat$ است. اما $D$ نمی‌تواند در فهرست باشد. زیرا بنا بر
ساخت، $n \in D$ اگر و تنها اگر $n\notin N_n$؛ ازاین‌رو $D \neq N_n$
برای هر $n \in \Nat$.
\end{proof}

\begin{explain}
برهان پیشین از قطر نام نبرد. بااین‌حال، اگر آن را چنین تصویر کنید،
می‌توانید آن را مشتمل بر یک قطر بدانید: فرض کنید مجموعه‌های $N_0$،
$N_1$، \dots، در آرایه‌ای نوشته شده‌اند که در آن $N_n$ را در سطر $n$اُم
با نوشتن $m$ در ستون $m$اُم، اگر و تنها اگر $m \in N_n$، نمایش می‌دهیم.
برای نمونه، فرض کنید چهار مجموعهٔ نخستِ آن فهرست
$\{0,1,2,\dots\}$، $\{1, 3, 5, \dots\}$، $\{0,1,4\}$ و
$\{2,3,4,\dots\}$ باشند؛ آنگاه آرایهٔ ما چنین آغاز می‌شود:
\[
\begin{array}{r@{}rrrrrrr}
  N_0 = \{ & \mathbf{0}, & 1, & 2, & & & & \dots\}\\
  N_1 = \{ &  & \mathbf{1}, &  & 3, &  & 5, & \dots\}\\
  N_2 = \{ & 0, & 1, &  &  & 4\phantom{,} &  & \}\\
  N_3 = \{ &  &  & 2, & \mathbf{3}, & 4, & & \dots\}\\
  &\vdots & & & & & & \ddots\phantom{\}}
  \end{array}
\]
آنگاه $D$ مجموعه‌ای است که با پایین‌رفتن روی قطر به دست می‌آید و $n
\in D$ را دقیقاً هنگامی برقرار می‌گیریم که $n$ روی قطر \emph{نباشد}.
پس در حالت بالا، $0$ و $1$ را کنار می‌گذاریم، $2$ را در مجموعه قرار
می‌دهیم، $3$ را کنار می‌گذاریم و الی آخر.
\end{explain}

\begin{prob}
با یک استدلال قطریِ صریح نشان دهید که مجموعهٔ همهٔ توابع
$f \colon \Nat \to \Nat$ !!{nonenumerable} است. یعنی نشان دهید که اگر
$f_1$، $f_2$، \dots، فهرستی از توابع باشد و هر $f_i\colon
\Nat \to \Nat$ باشد، آنگاه تابعی مانند $g \colon \Nat \to
\Nat$ وجود دارد که در این فهرست نیست.
\end{prob}

\end{document}
